\documentclass[12pt]{article}
\usepackage[a4paper]{geometry}
\usepackage[utf8]{inputenc}
\usepackage[english,main=russian]{babel}
\usepackage[T2A]{fontenc}
\usepackage{mathtools}
\usepackage{amsfonts}
\usepackage{minted}

\setminted{
  fontsize=\small,
  bgcolor=white,
  frame=none,
  breaklines=true
}
\title{Исчисление конструкций}
\begin{document}
Исчисление конструкций, которое является чисто функциональным языком и лежит в основе Исчисления индуктивных конструкций, было создано Тьерри Коканом (T. Coquand) и Жераром Юэ (G. Huet). Оно может быть определено как чистая система типов (PTS). Чистая система типов --- это типизированное лямбда-исчисление с синтаксисом, в котором единообразно описываются термы и типы. Термы включают в себя переменные, типизированные лямбда-абстракции (записываются как \mintinline{text}{fun x : A => t}) и применения (записываются как $t u$). Тип лямбда-абстракции --- это зависимое произведение $\Pi x : A, B$, дающее возможность типу $B$ зависеть от переменной $x$. Обозначение $A \rightarrow B$ --- это тип для функций из типа $A$ в тип $B$. Это частный случай $\Pi$-типа, когда тип $B$ не зависит от $x$.

Типы сами являются типизированными объектами, тип для типа --- это специальная константа, называемая \texttt{sort}. Существует по крайней мере один \texttt{sort}, называемый \texttt{Type}. Разные системы типов в своей основе имеют разные множества сортов: каждый сорт соответствует определенному универсуму объектов, также с каждым сортом можно делать свои произведения.

Например, если $A$ --- это тип, то тождественная функция \mintinline{text}{fun x : A => x} --- это терм, имеющий тип $A \rightarrow A$. Если $A$ --- это типовая переменная, мы можем сконструировать полиморфную тождественную функцию \texttt{fun A : Type => fun x : A => x}, имеющую тип $\Pi A : \texttt{Type}, A \rightarrow A$.

В исчислении конструкций у нас есть бесконечное множество сортов:

$$
S \stackrel{\text{def}}{=} \mathit{Prop} \cup \bigcup_{i \in \mathbb{N}} \mathit{Type}_i
$$

Сорт \texttt{Prop} --- это тип для выражений, которые обозначают логические высказывания. Мы будем следовать изоморфизму Карри -- Ховарда, где высказывание $A$ соответствует типу (типу доказательств $A$).

У нас есть язык программирования (типизированное $\lambda$-исчисление --- CIC, и логика высшего порядка). В соответствии с изоморфизмом Карри -- Ховарда каждому высказыванию в логике соответствует определенный тип в языке программирования. Каждый из этих типов имеет тип \texttt{Prop}. Доказательство высказывания --- это элемент типа, соответствующего данному высказыванию.

Если $A$ и $B$ --- два типа, соответствующие логическим высказываниям, то высказывание $A \implies B$ соответствует типу функций $A \rightarrow B$, преобразующих доказательства $A$ в доказательства $B$. Высказывание $A \land B$ представлено типом упорядоченных пар $A \times B$. Пусть дан тип $T$, тогда тип $\prod x : T, B(x)$ будет соответствовать высказыванию $\forall x : T, B(x)$. Поскольку типы представляют логические высказывания, язык будет содержать пустые типы, соответствующие недоказуемым высказываниям.

\section{Обозначения}

Мы будем использовать обозначение $\forall x : A, B(x)$ вместо $\prod x : A, B(x)$, когда $B$ представляет высказывание (имеет тип \texttt{Prop}). Мы будем писать $t[u/x]$ для терма $t$, в котором переменная $x$ заменяется термом $u$.

Исчисление индуктивных конструкций оперирует метавысказываниями (judgement-ами) вида:

$$
\underbrace{x_1 : A_1, \ldots , x_n : A_n}_{\text{контекст}} \vdash \underbrace{t : A}_{\text{Вывод}}
$$

В контексте переменная $x_i$ типа $A_i$ соответствует имени объекта. Если следовать парадигме <<propositions as types>>, то тип $A_i$ соответствует логическому высказыванию, а $x_i$ --- это имя гипотезы, что $A_i$ доказуемо.

Метавысказывание (judgement) можно прочитать следующим образом: при предположении, что если доказуемы высказывания $A_i$ ($x_i$ --- это имена для доказательств), то доказуемо высказывание $A$. Это логическая интерпретация метавысказывания. Интерпретация в рамках исчисления конструкций: если нам даны термы $x_i : A_i$, то мы можем построить правильный терм $t : A$.

\begin{minted}[breaklines, breakanywhere, tabsize=2]{coq}
Notation "A -> B" := (forall (_ : A), B)
                       (right associativity, at level 99).

Definition False_CoC : Prop :=
  forall P : Prop, P.

Definition not_CoC (P : Prop) : Prop :=
  P -> False_CoC.

Notation "~ A" := (not_CoC A) (at level 75, right associativity).

Definition And_CoC (A B : Prop) : Prop :=
  forall (C : Prop), (A -> B -> C) -> C.
\end{minted}


Из логики первого порядка возьмем термин ``сигнатура'' и построим формальную теорию.


\end{document}
